\subsection{8. Was der Operationssatz jetzt tatsächlich
hergibt}\label{was-der-operationssatz-jetzt-tatsuxe4chlich-hergibt}

{\def\LTcaptype{none} % do not increment counter
\begin{longtable}[]{@{}
  >{\raggedright\arraybackslash}p{(\linewidth - 4\tabcolsep) * \real{0.3333}}
  >{\raggedright\arraybackslash}p{(\linewidth - 4\tabcolsep) * \real{0.3333}}
  >{\raggedright\arraybackslash}p{(\linewidth - 4\tabcolsep) * \real{0.3333}}@{}}
\toprule\noalign{}
\begin{minipage}[b]{\linewidth}\raggedright
Vertrag
\end{minipage} & \begin{minipage}[b]{\linewidth}\raggedright
Mathematisch abgesichert
\end{minipage} & \begin{minipage}[b]{\linewidth}\raggedright
Nicht dadurch verfügbar
\end{minipage} \\
\midrule\noalign{}
\endhead
\bottomrule\noalign{}
\endlastfoot
Markierter endlicher Matrixcompiler & Matrizen, Ordnungen, konkrete
endliche Syntheseidentitäten & Vollständiges Fockinstrument,
Präparation, physischer Zeitgenerator \\
H allein & Modellzeitentwicklung & Unabhängiges Schalten von X und
\(N_b\) \\
X und \(N_b\), wenn als Kontrollen gewährt & N=3-Kontrollalgebra der
Dimension 14 & Beliebige Modenoperationen \\
Zusätzlich dokumentierter vorzeichenrichtiger Clock-Lift &
N=3-Kontrollalgebra der Dimension \textbf{84} & Vollständige
Zustandsunterscheidung oder Ladung-eins-Instrument \\
\end{longtable}
}

Die neue Clock-Rechnung wurde übernommen und reproduziert. Die fünf
Multiplizitätszeilen für die sechs Clockphasen sind:

{\def\LTcaptype{none} % do not increment counter
\begin{longtable}[]{@{}lrrrrrr@{}}
\toprule\noalign{}
Teilraum & 0 & 1 & 2 & 3 & 4 & 5 \\
\midrule\noalign{}
\endhead
\bottomrule\noalign{}
\endlastfoot
Dunkle F3 & 6384 & 6280 & 6280 & 6384 & 6280 & 6280 \\
Dunkle BF & 16 & 8 & 8 & 16 & 8 & 8 \\
Aktiv 7 & 560 & 440 & 440 & 560 & 440 & 440 \\
Aktiv 10 & 112 & 88 & 88 & 112 & 88 & 88 \\
Aktiv 12 & 80 & 40 & 40 & 80 & 40 & 40 \\
\end{longtable}
}

Die aktiven Zeilen tragen zusätzlich einen Zweiniveaufaktor. Alle 45504
N=3-Zustände sind erfasst; der Kommutant hat weiterhin Dimension
240742144. Die alte Zahl 14 beschreibt den engeren Zweikontrollvertrag,
nicht den jetzt zusätzlich geprüften Clockvertrag.

\subsubsection{8.1 Konkrete
Präparationsgrenze}\label{konkrete-pruxe4parationsgrenze}

Alle genannten Fockkontrollen erhalten N. Ein Wort aus ihnen und jeder
einzeln zahlenerhaltende ausgewählte Messzweig bleiben im
Ausgangssektor. Sie können aus dem leeren N=0-Zustand \textbf{nicht} den
N=64-Grundzustand erzeugen. Diese Aussage betrifft zahlenerhaltende
Krauszweige; eine bloß U(1)-kovariante offene Dynamik kann sehr wohl
geladene Krausoperatoren haben.

Auch aus dem voll besetzten F führt bloßes \(e^{-itH}\) nicht zur
Konvergenz auf \(\Omega\): Der Grundzustandsüberlapp bleibt konstant.
Normiertes \(e^{-\tau H}F\) konvergiert mathematisch mit der bewiesenen
Lücke, benötigt aber eine gesonderte Instrument- und
Ressourcenherleitung.

\subsubsection{8.2 Der kleinste konkrete
Ladungstest}\label{der-kleinste-konkrete-ladungstest}

Ein Kandidat ist eine explizite Referenzmode c mit
\(T=\lambda(c^\dagger f_r+f_r^\dagger c)\). Sie erhält die gemeinsame
Ladung, verändert aber die Ladung der ursprünglichen Bank. Der endliche
Austausch- und Detuningtest ist exakt; \textbf{dieses T wurde nicht als
neuer nativer Hamiltonterm eingeführt}.

Für dieselbe N=64-Präparation unter \(H+\mu N\) gilt
\(G_\mu(z)=G_0(z-\mu)\), insbesondere \(m_1=\mu\). Das misst nur gegen
eine festgelegte Referenz: Ein gemeinsamer Zusatz
\(\mu(N_{\rm Bank}+N_{\rm Referenz})\) bleibt unsichtbar. Die Spektral-
schranken dieses Berichts gelten als Grundzustandsschranken nur für μ=0.

\subsection{9. Relativistisches Wörterbuch: Ausschluss und zwei präzise
Alternativen}\label{relativistisches-wuxf6rterbuch-ausschluss-und-zwei-pruxe4zise-alternativen}

Die Indexfrage darf nicht nachträglich durch einen Namen beantwortet
werden. Die Konventionen für Zweikomponentenspinoren wurden mit
\href{https://arxiv.org/abs/0812.1594}{Dreiner, Haber und Martin}
abgeglichen; die folgenden speziellen W-Tests sind eigene exakte
Tensorrechnungen.

\subsubsection{9.1 Was nicht funktioniert}\label{was-nicht-funktioniert}

Für 64 gleichhändige Weylfelder mit innerem Index I und einen skalaren
Vermittler ist \(M_A\otimes\varepsilon_{\rm Lorentz}\) symmetrisch. Die
Grassmann-Antisymmetrisierung verschwindet daher identisch, für alle 60
W-Zeilen. Außerdem lässt die volle irreduzible innere Darstellung
\(16\otimes4\) auf denselben 64 Komponenten nach Schur keine zusätzliche
kommutierende nichttriviale Lorentzspinorwirkung zu.

\subsubsection{9.2 Bereits bekannter nichtverschwindender
Typ}\label{bereits-bekannter-nichtverschwindender-typ}

Ein symmetrischer Lorentzspinortensor koppelt an das antisymmetrische M.
Der Vermittler hat dann einen passenden dualen (1,0)-Typ samt
adjungiertem Typ. Diese Variante besteht den algebraischen
Nichtnulltest, aber noch keinen vollständigen Kinetik-, Positivitäts-,
Constraints- oder Herkunftstest.

\subsubsection{9.3 Neue minimale skalare Alternative - ausdrücklich
zusätzlich}\label{neue-minimale-skalare-alternative---ausdruxfccklich-zusuxe4tzlich}

Will man zugleich W, gleiche Weylhand, lokale Bilinearität und einen
skalaren Vermittler behalten, kann man einen \textbf{unabhängigen}
Hilfsindex a=1,2 mit alternierender Form einführen:

\[\adjustbox{max width=\linewidth}{$\displaystyle 
b_A^\dagger(M_A)_{IJ}\varepsilon_{ab}
\varepsilon_{\alpha\beta}\psi_{I a\alpha}\psi_{J b\beta}
+\mathrm{h.c.}
$}\]

M und \(\varepsilon_{ab}\) sind jeweils antisymmetrisch; ihr Produkt ist
als innerer Kopplungstensor symmetrisch. Zusammen mit der Lorentz-
Epsilonform ist der Gesamttensor antisymmetrisch und nicht null. Alle 60
Kanäle bestehen den Test. Eine nichtverschwindende alternierende Form
existiert nicht in Dimension eins; in Dimension zwei ist sie bis auf
Normierung eindeutig. \textbf{In dieser eng benannten Klasse von
Tensorfaktor-Reparaturen ist die binäre Ergänzung minimal.}

Das ist noch keine gefundene native Lösung: Sie verdoppelt die inneren
Weylkomponenten von 64 auf 128, zusätzlich zu deren Lorentzspinorindex.
Ein Double-Cover-Minuszeichen stellt nicht automatisch zwei unabhängige
Felder bereit. Auch Nambu-Umbenennung \((f,f^\dagger)\) genügt nicht:
Das gemischte Produkt hat Ladung null statt −2, sodass derselbe
\(b^\dagger ff\)-Ladungsvertrag nicht erhalten bleibt.

Der entscheidende Quellenauftrag ist daher eng: Gibt es diese zweite
gleichgeladene, CAR-unabhängige Komponente bereits im tatsächlichen
Compilerprozess? Und liefert ihre Projektion genau das bisherige H und
die geprüfte Antwort? Ohne beides bleibt die skalare Alternative ein
zusätzliches Modell, nicht eine Erklärung des ursprünglichen.

\subsection{10. Nächste Schritte mit eindeutiger
Erfolgskontrolle}\label{nuxe4chste-schritte-mit-eindeutiger-erfolgskontrolle}

\begin{enumerate}
\def\labelenumi{\arabic{enumi}.}
\tightlist
\item
  \textbf{Native Quelle des Austauschoperators bestimmen.} Ein
  tatsächliches Operationswort mit Anfangszustand, Detektor,
  Adjungiertem, Ladungsbilanz und Record angeben. Ein weiteres neutrales
  Wort oder bloßer Algebraabschluss schließt diese Aufgabe nicht.
\item
  \textbf{Den Rest der nativen Antwort kontrollieren.} Das nächste Ziel
  ist \(\Sigma_2\) beziehungsweise sein erster Norm- und
  Momentkoeffizient, gemeinsam in N=63,64,65. Die vorliegenden
  Summenregeln und Polschranken sind zwingende Akzeptanztests. Den Rest
  auf null zu setzen ist exakt ausgeschlossen; eine Näherung braucht
  eine Restfehlergrenze.
\item
  \textbf{Die beiden Feldtypen an der Quelle entscheiden.} Entweder der
  symmetrische Vermittler mit korrektem Adjungierten und Kinetik, oder
  der skalare Typ mit wirklich nachgewiesener zweiter Komponente. Erst
  Nichtnullkopplung, Symmetrie, Ladung, CAR und positive Kinetik
  gemeinsam zählen als Feldadapter.
\item
  \textbf{Erst danach zwei operational bestimmte Teile verbinden.} Für
  einen tatsächlich hergeleiteten ungeraden Austausch wäre das
  projizierte Ein-Loch-Transfermatrixelement proportional zum jetzt
  eingeschlossenen Residuum. Ein solches formales Matrixelement beweist
  weder Verfügbarkeit des Austauschs noch einen isolierten
  Zweibank-Gesamtraum oder höhere Störungsordnungen. Eine räumliche
  Skalierung wurde hier nicht vorgezogen.
\end{enumerate}

{\def\LTcaptype{none} % do not increment counter
\begin{longtable}[]{@{}
  >{\raggedright\arraybackslash}p{(\linewidth - 4\tabcolsep) * \real{0.3333}}
  >{\raggedright\arraybackslash}p{(\linewidth - 4\tabcolsep) * \real{0.3333}}
  >{\raggedright\arraybackslash}p{(\linewidth - 4\tabcolsep) * \real{0.3333}}@{}}
\toprule\noalign{}
\begin{minipage}[b]{\linewidth}\raggedright
Tor
\end{minipage} & \begin{minipage}[b]{\linewidth}\raggedright
Nutzen dieser Revision
\end{minipage} & \begin{minipage}[b]{\linewidth}\raggedright
Entscheidender verbleibender Nachweis
\end{minipage} \\
\midrule\noalign{}
\endhead
\bottomrule\noalign{}
\endlastfoot
T1 & Exakter Clock-Kontrollabschluss und engere Ressourcenfrage &
Ursprüngliches vollständiges Operations- und Rahmenwörterbuch \\
T2 & Geladene native Antwort, sichtbarer Pol, drittes Moment &
Quelleneinbettung und renormiertes Half-Charge-Feld mit
Energie/Adjungiertem \\
T3 & Präzise Anforderung an ungeraden Austausch & Gemeinsamer
operationaler räumlicher, schließlich 3+1D-Träger \\
T4 & Falscher Skalartyp ausgeschlossen, minimale Alternativen & Chirales
Maß, Anomalien, Spiegelkontrolle, vollständige Feldkinetik \\
T5 & Pole und Reststruktur lokal kontrolliert & Gemeinsamer
wechselwirkender Grenzwert, Clusterstruktur, Streuung \\
T6 & Strengere interne/Lorentz-Indexbilanz & Familien, Massen und
Kopplungen auf demselben physikalischen Träger \\
T7 & Keine neue Spin-2-Konstruktion & Dynamischer Spin 2, Helizitäten
und universelle Kopplung \\
T8 & Eindeutiger Modellgrundzustand; Präparationslücke konkret &
Primitive Zustandswahl, Ressourcen, Records und Instrumente \\
\end{longtable}
}

\subsection{11. Erratum und
Versionsdisziplin}\label{erratum-und-versionsdisziplin}

Im eigenen v1.6.2-Text fehlte in Abschnitt 5.2 zwischen Entnahme- und
Additionsresolvente ein Pluszeichen. Richtig ist auf der dortigen
N=4-Referenz \(G=Z_h/(z+\epsilon_h)+a^\dagger(z+E_--H_5)^{-1}a\). Die
früheren Prüfrechnungen verwendeten die additive CAR-Gewichtsregel; der
Darstellungsfehler wird hier ausdrücklich korrigiert. Die archivierte
v1.6.2 wird nicht stillschweigend umgeschrieben.

Die zugelieferte v1.6.3 bleibt ebenfalls unverändert. Ihre Pol- und
Clockbefunde sind als übernommene und frisch reproduzierte Ergebnisse
kenntlich. Die engeren Schranken, das dritte Moment, der Ausschluss der
exakten Zweilinienantwort und die minimale skalare Hilfsindexalternative
sind die eigenen zusätzlichen Ergebnisse dieser Konsolidierung. Kein
literaturweiter Neuheitsanspruch, keine experimentelle Bestätigung und
kein Gesamtabschluss werden behauptet.

\begin{center}\rule{0.5\linewidth}{0.5pt}\end{center}

